\documentclass[10pt,a4paper]{article}

\usepackage[utf8]{inputenc}
\usepackage[T1]{fontenc}
\usepackage[french]{babel}
\usepackage{lmodern}
\usepackage[margin=1.35cm]{geometry}
\usepackage{xcolor}
\usepackage{booktabs}
\usepackage{array}
\usepackage{pgfplots}
\usepackage{titlesec}
\usepackage{microtype}

\definecolor{bleu}{HTML}{1557A0}
\definecolor{bleuclair}{HTML}{EAF2FB}
\pgfplotsset{compat=1.18}
\setlength{\parindent}{0pt}
\setlength{\parskip}{3pt}
\titleformat{\section}{\large\bfseries\color{bleu}}{}{0pt}{}
\titleformat{\subsection}{\normalsize\bfseries\color{bleu}}{}{0pt}{}
\titlespacing*{\section}{0pt}{7pt}{3pt}
\titlespacing*{\subsection}{0pt}{5pt}{2pt}

\begin{document}

\begin{center}
    {\LARGE\bfseries\color{bleu} Benchmark d'un arbre de recherche binaire}\\[-1pt]
    {\large Insertion, recherche et suppression}\\[2pt]
    {\normalsize ERRAMI Youssef \quad -- \quad C++}
\end{center}
\vspace{-4pt}
\hrule

\section{Objectif}
Mesurer le temps total nécessaire pour insérer, rechercher et supprimer
$n$ valeurs dans un arbre de recherche binaire (ABR), puis comparer les
résultats pour différentes tailles d'entrée.

\section{Configuration matérielle}
\renewcommand{\arraystretch}{1.05}
\begin{tabular}{@{}>{\bfseries}p{3.1cm}p{12.2cm}@{}}
Ordinateur & MSI Stealth GS66 12UGS\\
Processeur & Intel Core i7-12700H, 14 cœurs\\
Carte graphique & NVIDIA GeForce RTX 3070 Ti\\
Mémoire vive & 16 Go DDR5\\
\end{tabular}

\section{Méthodologie}
Pour chaque taille $n$, le programme crée les entiers de $0$ à $n-1$,
puis les mélange avant de les insérer. Il chronomètre successivement
l'insertion des $n$ valeurs, la recherche de ces mêmes valeurs (recherches
réussies) et leur suppression. Les durées affichées sont des temps totaux
en millisecondes. L'arbre du test est temporaire et son affichage trié est
enregistré dans \texttt{output.txt}.

\textbf{\color{bleu}Protocole conseillé :} exécuter le test pour
$n=1\,000$, $5\,000$, $10\,000$, $50\,000$ et $100\,000$. Pour chaque
taille, reporter les trois temps affichés dans les graphiques ci-dessous.
Les zéros sont des champs à remplacer par les mesures réelles ; aucun
résultat n'est inventé dans ce rapport.

\section{Résultats expérimentaux}
\begin{center}
\small
\begin{minipage}{0.32\textwidth}
\centering\textbf{Insertion}\\[-2pt]
\begin{tikzpicture}
\begin{axis}[
    width=\linewidth,height=4.5cm,
    xlabel={Nombre de valeurs},ylabel={Temps (ms)},
    symbolic x coords={1k,5k,10k,50k,100k}, xtick=data,
    x tick label style={font=\scriptsize,rotate=35,anchor=east},
    y tick label style={font=\scriptsize},
    label style={font=\scriptsize}, grid=major,
    grid style={bleuclair}, tick style={bleu},
    every axis plot/.append style={very thick,bleu,mark=*}
]
\addplot coordinates {(1k,0) (5k,0) (10k,0) (50k,0) (100k,0)};
\end{axis}
\end{tikzpicture}
\end{minipage}\hfill
\begin{minipage}{0.32\textwidth}
\centering\textbf{Recherche}\\[-2pt]
\begin{tikzpicture}
\begin{axis}[
    width=\linewidth,height=4.5cm,
    xlabel={Nombre de valeurs},ylabel={Temps (ms)},
    symbolic x coords={1k,5k,10k,50k,100k}, xtick=data,
    x tick label style={font=\scriptsize,rotate=35,anchor=east},
    y tick label style={font=\scriptsize},
    label style={font=\scriptsize}, grid=major,
    grid style={bleuclair}, tick style={bleu},
    every axis plot/.append style={very thick,bleu,mark=square*}
]
\addplot coordinates {(1k,0) (5k,0) (10k,0) (50k,0) (100k,0)};
\end{axis}
\end{tikzpicture}
\end{minipage}\hfill
\begin{minipage}{0.32\textwidth}
\centering\textbf{Suppression}\\[-2pt]
\begin{tikzpicture}
\begin{axis}[
    width=\linewidth,height=4.5cm,
    xlabel={Nombre de valeurs},ylabel={Temps (ms)},
    symbolic x coords={1k,5k,10k,50k,100k}, xtick=data,
    x tick label style={font=\scriptsize,rotate=35,anchor=east},
    y tick label style={font=\scriptsize},
    label style={font=\scriptsize}, grid=major,
    grid style=bleuclair, tick style={bleu},
    every axis plot/.append style={very thick,bleu,mark=triangle*}
]
\addplot coordinates {(1k,0) (5k,0) (10k,0) (50k,0) (100k,0)};
\end{axis}
\end{tikzpicture}
\end{minipage}
\end{center}

\section{Interprétation}
Compléter les graphiques avec les mesures obtenues, puis comparer
l'évolution des trois temps lorsque $n$ augmente. La hauteur de l'ABR
influence directement les opérations : un arbre équilibré donne en moyenne
des parcours de l'ordre de $\log(n)$, tandis qu'un arbre très déséquilibré
peut nécessiter jusqu'à $n$ étapes. Les valeurs sont mélangées dans ce test,
mais le résultat dépend aussi de l'ordre aléatoire généré.

\vspace{2pt}
\colorbox{bleuclair}{\parbox{\dimexpr\textwidth-2\fboxsep\relax}{%
\textbf{\color{bleu}Compilation :} \texttt{g++ -std=c++17 -O2 arbre\_recherche\_binaire.cpp -o arbre.exe}
\quad
\textbf{\color{bleu}Exécution :} \texttt{.\textbackslash arbre.exe}}}

\end{document}
